\chapter{引言}
\section{研究背景}
时间序列预测可用于描述工程系统的短期变化。建立预测模型时，应说明观测变量、采样间隔和适用范围，并在独立样本上报告误差。Box 等\cite{box2015}系统讨论了时间序列的分析与预测方法。

\section{研究内容}
本文的示例内容用于检查目录、页眉和正文样式。图、表及公式的编号均可使用交叉引用自动更新。

\subsection{观测数据}
本例以等间隔采样序列为对象，采样间隔设为 \qty{1}{\second}。观测变量为位移，单位采用 \unit{\milli\metre}。

% 这里的 \clearpage 用于演示偶数页页眉。写论文时可删除。
\clearpage
\section{论文结构}
第 2 章演示公式、插图和三线表的排版。第 3 章给出结论的写作示例。文末列出正文引用的文献，并在附录中给出计算示例。
